\section{Minimum orbits from root and singular-value data}\label{sec:minima41}
The width theorem becomes a numerical exponent only after its orbit
invariant is evaluated. This section evaluates it on every compact simple
adjoint representation, on the defining classical representations, and
on an infinite family of fundamental representations. The root and
normal-form tools are classical; they are not asserted to originate in
this paper. What they determine here is the dynamic width exponent under
the action and calibration hypotheses of
Theorem~\ref{thm:intrinsic-width41}. We use the conventional root systems
and highest-weight notation of \cite[Sections 21--23, 27, 30--32]{Etingof}.

\subsection{All compact simple adjoint representations}
Let $\mathfrak h$ be a compact simple real Lie algebra, with its invariant
positive inner product, and let $H$ be its connected adjoint group.
Choose a maximal torus and a simple system $\Delta$ in the root system
$\Phi$ of its complexification. For $I\subset\Delta$ let $\Phi_I$ be the
root subsystem generated by $I$; the empty subsystem has zero roots.

\begin{theorem}[Adjoint minimum-orbit formula]\label{thm:adjoint41}
For the adjoint action on $\mathfrak h$,
\begin{equation}\label{eq:adjoint41}
 s(\operatorname{Ad})=|\Phi|-
       \max_{\alpha\in\Delta}|\Phi_{\Delta\setminus\{\alpha\}}|.
\end{equation}
Every minimizing unit vector is conjugate to a point on a fundamental
coweight ray corresponding to a maximizing deletion in
\eqref{eq:adjoint41}. Conversely every such ray minimizes. The compact
centralizer has Lie algebra consisting of the deleted-node Levi algebra
and a one-dimensional center. The values are listed in
Table~\ref{tab:adjoint41}.
\end{theorem}
\begin{proof}
Every element is conjugate into a closed Weyl chamber of the compact
Cartan algebra. We identify that algebra with its real root-coordinate
space. If $X$ lies in this chamber, put
$I(X)=\{\alpha\in\Delta:\alpha(X)=0\}$. A positive root is a
nonnegative integral combination of simple roots. Consequently its
value on $X$ vanishes exactly when it belongs to $\Phi_{I(X)}$.
The real centralizer dimension is therefore
$\operatorname{rank}\Phi+|\Phi_{I(X)}|$, while
$\dim\mathfrak h=\operatorname{rank}\Phi+|\Phi|$. Hence
\[
 \dim(HX)=|\Phi|-|\Phi_{I(X)}|.
\]
For $X\ne0$, the simple roots span the dual Cartan space, so $I(X)$
is proper. Every proper subset is contained in a subset obtained by
one deletion. If it is strictly contained, adjoining a missing simple
root adds at least that root and its negative, and strictly increases
the root count. Thus the largest possible centralizer occurs precisely
at a maximizing single deletion. Points of the closed chamber with
exactly that zero set are the positive fundamental coweight ray.
Normalization chooses its unit point. This proves the formula and all
the assertions about the minimizers and their centralizers.

For completeness the root counts entering the table can be obtained
without an orbit-type classification. For $A_\ell$, deletion of node
$i$ leaves $A_{i-1}+A_{\ell-i}$ and has
$i(i-1)+(\ell-i)(\ell-i+1)$ roots. For $B_\ell,C_\ell$ the analogous
count is $i(i-1)+2(\ell-i)^2$. These convex quadratic expressions are
maximized at endpoints; they give the displayed counts, including the
rank-two ties. For $D_\ell$ the chain deletions give
$i(i-1)+2(\ell-i)(\ell-i-1)$, while a spin-end deletion gives
$\ell(\ell-1)$. Their maximum is $2(\ell-1)(\ell-2)$ for
$\ell\ge4$, with the familiar ties at $D_4$.
The exceptional diagrams have at most eight nodes: deleting them gives
respectively a maximal $A_1$, $B_3$ or $C_3$, $D_5$, $E_6$, and $E_7$
root subsystem. Their root counts are $2,18,40,72,126$, whereas the
original systems have $12,48,72,126,240$ roots. These counts follow by
reflection from their simple roots, or directly from the standard
coordinate root lists. Substitution completes the table.
\end{proof}

\begin{table}[ht]
\centering
\caption{Compact adjoint orbits and their conditional width exponents.
The central one-dimensional torus in the centralizer is not a root.}
\label{tab:adjoint41}
\begin{tabular}{llll}
\toprule
Type & Largest deleted root subsystem & $s(\operatorname{Ad})$ & $s/2$\\
\midrule
$A_\ell\ (\ell\ge1)$ & $A_{\ell-1}$ & $2\ell$ & $\ell$\\
$B_\ell\ (\ell\ge2)$ & $B_{\ell-1}$ & $4\ell-2$ & $2\ell-1$\\
$C_\ell\ (\ell\ge2)$ & $C_{\ell-1}$ & $4\ell-2$ & $2\ell-1$\\
$D_\ell\ (\ell\ge4)$ & $D_{\ell-1}$ & $4\ell-4$ & $2\ell-2$\\
$G_2$ & $A_1$ & $10$ & $5$\\
$F_4$ & $B_3$ or $C_3$ & $30$ & $15$\\
$E_6$ & $D_5$ & $32$ & $16$\\
$E_7$ & $E_6$ & $54$ & $27$\\
$E_8$ & $E_7$ & $114$ & $57$\\
\bottomrule
\end{tabular}
\end{table}
These are orbits of actual real vectors in the compact Lie algebra.
They are not complex nilpotent orbits or projective highest-weight
orbits. Confusing these objects changes the dimensions. The table is
an evaluation of a classical centralizer calculation, not a claim of
new classification of compact simple Lie algebras.

\subsection{Defining and low-rank spin representations}
\begin{proposition}[Transitive defining actions]\label{prop:standard41}
The defining real actions of $\SO(n)$ on $\R^n$, $\SU(n)$ on the
underlying real space of $\C^n$, and $\operatorname{Sp}(n)$ on
$\mathbb H^n$ have minimum orbit dimensions
\[
 n-1,\qquad 2n-1,\qquad 4n-1,
\]
respectively ($n\ge3$, $n\ge2$, and $n\ge1$). Their unit-vector
stabilizers are $\SO(n-1)$, $\SU(n-1)$, and
$\operatorname{Sp}(n-1)$. The real spin representations of
$\operatorname{Spin}(3)$ and $\operatorname{Spin}(5)$, and either
complex half-spin representation of $\operatorname{Spin}(6)$ viewed
as a real representation, have minimum orbit dimensions $3,7,7$.
\end{proposition}
\begin{proof}
Extend a unit vector to an orthonormal basis over the respective
field. In the special orthogonal and special unitary cases, adjust
the determinant on its orthogonal complement. This proves transitivity
on the sphere. A matrix fixing the first basis vector acts freely by
the indicated smaller group on its complement, giving the stabilizers
and the dimensions.
The standard low-rank identifications
$\operatorname{Spin}(3)=\operatorname{Sp}(1)$,
$\operatorname{Spin}(5)=\operatorname{Sp}(2)$ and
$\operatorname{Spin}(6)=\SU(4)$ identify the stated spin modules with
these defining modules (the two half-spin modules in the last case are
dual). They follow also by matching the $B_1$, $B_2=C_2$, and
$D_3=A_3$ fundamental weights and dimensions
\cite[Sections 30--32]{Etingof}. Apply the preceding transitivity.
\end{proof}
This proposition does not assert a formula for all higher spin modules.
The adjoint table, including every exceptional type, and the next
fundamental family have separate complete proofs.

\subsection{A fundamental family whose minimum need not be a highest-weight orbit}
Let $\SU(n)$ act on complex skew-symmetric matrices by
$A\mapsto gAg^T$. This is $\bigwedge^2\C^n$, with its underlying
real Euclidean structure. A fixed scalar normalization of the
Hilbert--Schmidt norm does not affect orbit dimensions.

\begin{theorem}[Second exterior powers]\label{thm:exterior41}
For this action on its real unit sphere,
\begin{equation}\label{eq:exterior41}
 s\bigl((\bigwedge^2\C^n)_{\R}\bigr)=
 \begin{cases}
 5,&n=3,4,\\
 13,&n=5,\\
 14,&n=6,\\
 4n-7,&n\ge7.
 \end{cases}
\end{equation}
For $n=3$ and $n\ge5$ the underlying real representation is irreducible.
For $n=4$ it is the sum of two equivalent real six-dimensional
irreducibles, each with minimum orbit dimension five. At $n=6$ a
minimizing vector is a full-rank symplectic form, with stabilizer
$\operatorname{Sp}(3)$ and orbit $\SU(6)/\operatorname{Sp}(3)$;
the decomposable highest-weight vector instead has orbit dimension
seventeen.
\end{theorem}
\begin{proof}
We recall the skew-symmetric unitary congruence normal form and its
argument. On $\C^n$ define the antilinear map $Fv=A\bar v$.
Skew symmetry implies $F^2=-AA^*$. On a positive eigenspace of
$AA^*$ with eigenvalue $a^2$, $F/a$ is antiunitary, squares to $-I$,
and pairs a unit vector with an orthogonal unit vector. Inductively
choosing such pairs gives blocks $aJ$, where
$J=\left(\begin{smallmatrix}0&1\\-1&0\end{smallmatrix}\right)$,
and a zero block. This is the classical unitary congruence construction
\cite{Youla}; the pairing argument supplies the needed special case.
Although the basis change can have determinant different from one,
conjugation by it preserves $\SU(n)$ and therefore the dimension of
the stabilizer inside $\SU(n)$.

Suppose the distinct positive singular values have multiplicities
$2k_1,\ldots,2k_b$, with $k=\sum_i k_i$ and $l=n-2k$.
A stabilizer preserves the eigenspaces of $AA^*$ and commutes there
with the associated quaternionic structure. Its unitary factors are
therefore $\operatorname{Sp}(k_i)$; on the kernel they form $U(l)$.
The symplectic factors have determinant one, so imposing determinant
one leaves $\SU(l)$ on the kernel when $l>0$. The stabilizer dimension
is
\[
 \sum_i k_i(2k_i+1)+
 \begin{cases}l^2-1,&l>0,\\0,&l=0.\end{cases}
\]
At fixed rank, coalescing the positive singular values increases this
quantity, because merging parts $a,b$ increases it by $4ab$.
Thus the smallest orbit at rank $2k<n$ has dimension
\begin{equation}\label{eq:skew-rank41}
 n^2-1-k(2k+1)-((n-2k)^2-1)
                         =k(4n-6k-1).
\end{equation}
At full rank $n=2k$ it has dimension
\begin{equation}\label{eq:skew-full41}
                 n^2-1-k(2k+1)=n(n-1)/2-1.
\end{equation}
The quadratic in \eqref{eq:skew-rank41} is concave in $k$, so its
minimum is at an endpoint. Comparing those endpoints with
\eqref{eq:skew-full41}, when present, gives exactly
\eqref{eq:exterior41}. These minimizing strata are attained by
normalized block forms, so this is not merely a dimension lower bound.

The complex module has highest weight $\omega_2$ and dual weight
$\omega_{n-2}$; for $n\ne4$ these are different when $n\ge3$.
Consequently its underlying real module is irreducible of complex type.
For $n=4$ the conjugate-linear Hodge involution on two-forms has square
one and commutes with $\SU(4)$; its fixed real six-space and its
imaginary multiple give the two equivalent real forms. The identification
$\SU(4)=\operatorname{Spin}(6)$ makes each the standard real
six-dimensional module, so its minimum orbit dimension is five.
Alternatively the full-rank equality in \eqref{eq:skew-full41} gives
the same whole-space minimum. Finally $\dim\SU(6)-\dim\operatorname{Sp}(3)
=35-21=14$, whereas \eqref{eq:skew-rank41} at $k=1$ gives seventeen.
\end{proof}

Under the intrinsic action-gap and calibration hypotheses, divide the
numbers in \eqref{eq:exterior41} by two to obtain the width exponents.
In particular the $\SU(6)$ fundamental module has order $N^7$, not
$N^{17/2}$; replacing the minimizing real-vector orbit by the
highest-weight orbit would give the wrong answer. These results do not
purport to classify every highest-weight representation. They compute
the invariant, and its minimizing stabilizers, on the stated complete
families rather than leaving those inputs as an existential rank test.
